%
% Settings file
%
%=======================================================================
%-------------This file controls various document settings--------------
%=======================================================================

%
% Should we compile for the Iliad (ebook reader), or regular A4
%
\newif\ifcompileiliad 
	\compileiliadfalse %Iliad (true) or regular A4 (false)?


%
% Set the fontsize and baselineskip here, use 13/15 for final document downscaling to 80% (UiB standard)
%
\newcommand{\TextSize}{13}
\newcommand{\BaseLineSkip}{15}
\newif\ifDownscaledFinalDoc
	\DownscaledFinalDoctrue		% Uib 13pt (true) or Regular 12pt (false)


%
% Compile draft (double line spacing)?
%
\newif\ifDraft
	\Draftfalse		% Draft (true) or Final (false)
	% Document settings file

%
% We use the book class
%
\ifcompileiliad
	\documentclass[10pt]{book}
\else
	\documentclass[12pt]{book}
\fi

%
% The geometry package allows consistent control of page layout
%
\ifcompileiliad
	\setlength{\parindent}{0pt}
	\setlength{\parskip}{3ex}
	\usepackage[papersize={122mm,160mm},
				hmargin=0.2cm,
				vmargin=0.2cm,
				tmargin=1.0cm,
				bmargin=0.4cm,
				dvips]
				{geometry}
\else
	\usepackage[paper = a4paper,         %To be shrunk to 80% when printed
				twoside,                 %Two-side mode, switches margins on 
				bindingoffset = 2mm,     %Offset for binding side of page
				hmargin = 25mm,          %Left and right margin
				vmargin = 25mm,          %Top and bottom margin
				dvipdfm]
				{geometry}
\fi	

%
% Font settings
%
\usepackage[T1]{fontenc}	% Activate Type 1 fonts
\usepackage{mathptmx}       % Use Times font, also for math
\usepackage[scaled]{helvet} % for sans serif fonts (\textsf{...} or \sffamiliy)
\usepackage{sectsty}        % Change section and chapter header
\allsectionsfont{\usefont{OT1}{phv}{bc}{n}\selectfont} % Set chapter/section header to narrow helvectica (arial-like)
%\usepackage[scaled]{luximono} % for monospaced fonts (\texttt{...} or \ttfamily)

% Fourier for math | Utopia (scaled) for rm | Helvetica for ss | Latin Modern for tt
%\usepackage{fourier} % math & rm
%\usepackage[scaled=0.875]{helvet} % ss
%\renewcommand{\ttdefault}{lmtt} %tt

%% Latin Modern (similar to CM with more characters)
%\usepackage{lmodern} % math, rm, ss, tt
%\usepackage[T1]{fontenc}
%
%% Palatino for rm and math | Helvetica for ss | Courier for tt
%\usepackage{mathpazo} % math & rm
%\linespread{1.05}        % Palatino needs more leading (space between lines)
%\usepackage[scaled]{helvet} % ss
%\usepackage{courier} % tt
%\normalfont
%\usepackage[T1]{fontenc}
%
%% Euler for math | Palatino for rm | Helvetica for ss | Courier for tt
%\renewcommand{\rmdefault}{ppl} % rm
%\linespread{1.05}        % Palatino needs more leading
%\usepackage[scaled]{helvet} % ss
%\usepackage{courier} % tt
%\usepackage{euler} % math
%%\usepackage{eulervm} % a better implementation of the euler package (not in gwTeX)
%\normalfont
%\usepackage[T1]{fontenc}
%
%% Times for rm and math | Helvetica for ss | Courier for tt
%\usepackage{mathptmx} % rm & math
%\usepackage[scaled=0.90]{helvet} % ss
%\usepackage{courier} % tt
%\normalfont
%\usepackage[T1]{fontenc}




%
% Packages we need
%
\usepackage{graphicx}       % Handles figures
\usepackage[utf8]{inputenc} % We want æøå
\usepackage{amsmath}
\usepackage{amsfonts}
\usepackage{amssymb}
\usepackage{amsthm}
\usepackage[mathcal]{euscript} % For caligraphy fonts
\usepackage{booktabs}       % Publication quality tables
\usepackage{setspace}       % Easy setting of line spacing
\usepackage{forloop}		% For-loops!
\usepackage[numbers,
			sort&compress
			]{natbib}       % Sort numerical keys for multiple cites
\usepackage{hypernat}       % Avoid breaking 'backref' option of hyperref package
\usepackage{float}
\usepackage[dvipsnames]{color}   % named colors
\usepackage[procnames]{listings} % beautiful listings
\usepackage{units}				 % semantically represent numbers with units

%
% The hyperref package allows advanced hyperlinking functionality
%
\usepackage[dvipdfmx,        % Use dvipdf driver
			backref,         % List citing occurences in the References
			colorlinks,      % Colored links
			citecolor=blue,  % Color of cite links
			linkcolor=blue,  % Color of links
			urlcolor=blue,   % Color of urls
			]{hyperref}
%\newcommand{\aref}[1]{\autoref*{#1}} % Prevent links
\newcommand{\aref}[1]{\autoref{#1}} % Enable links


%
% Set fancy page header using fancyhdr package
% 
\usepackage{fancyhdr}
\pagestyle{fancy}

% Ensure that the chapter and section headings are in lowercase
\setlength{\headheight}{15pt}
\renewcommand{\chaptermark}[1]{\markboth{#1}{}}
\renewcommand{\sectionmark}[1]{\markright{\thesection\ #1}}

% Delete current section for header/footer
\fancyhf{}

% Define header/footer layout
\fancyhead[LE,RO]{\bfseries\thepage}
\fancyhead[LO]{\bfseries\rightmark}
\fancyhead[RE]{\bfseries\leftmark}
\renewcommand{\headrulewidth}{0.5pt}
%%%%%%%%%%%%%%%%%%%%%%%%%%%%%%%%%%%%%%%%%%%%%%%%%%%%%%%%%%%%%%%%%%%%%%%%%%%%%%%%%%%%%%
%
%
\newcommand{\Sfrac}[2]{{\textstyle{\frac{#1}{#2}}}}
\newcommand{\ddd}[1]{{\frac{\partial}{\partial #1}}}

\newcommand{\ohf}{\frac{1}{2}}
\newcommand{\othrd}{\frac{1}{3}}

\newcommand{\scrO}{{\mathcal{O}}}

\newcommand{\R}{\mathbb{R}}

\newcommand{\tx}{\tilde{x}}
\newcommand{\ty}{\tilde{y}}
\newcommand{\tz}{\tilde{z}}
\newcommand{\ttt}{\tilde{t}}
\newcommand{\teta}{\tilde{\eta}}
\newcommand{\tPhi}{\tilde{\Phi}}
\newcommand{\tD}{\tilde{D}}
\newcommand{\tH}{\tilde{H}}
\newcommand{\tu}{\tilde{u}}
\newcommand{\tr}{\tilde{r}}
\newcommand{\ts}{\tilde{s}}
\newcommand{\tK}{\tilde{K}}

\newcommand{\ta}{\alpha}
\newcommand{\tambda}{\lambda}
\newcommand{\sca}{-\kappa/\mu^\nu}

\newcommand{\tsrK}{\tilde{\mathcal{K}}}
\newcommand{\srK}{\mathcal{K}}
\newcommand{\srF}{\mathcal{F}}

\DeclareMathOperator{\sech}{sech}

\newcommand*\Lop{\mathcal{L}}
\newcommand*\Aop{\mathcal{A}}
\newcommand*\RR{\mathbf{R}}
\newcommand*\NN{\mathbf{N}}
\newcommand*\ZZ{\mathbf{Z}}
%%%%%%%%%%%%%%%%%%%%%%%%%%%%%%%%%%%%%%%%%%%%%%%%%%%%%%%%%%%%%%%%%%%%%%%%%%%%%
%
%
%
\renewcommand{\theequation}{\arabic{section}.\arabic{equation}}
\theoremstyle{plain}
\newtheorem{theorem}{\bf Theorem}[section]
\newtheorem{lemma}[theorem]{\bf Lemma}
\newtheorem{corollary}[theorem]{\bf Corollary}
\newtheorem{proposition}[theorem]{\bf Proposition}
%\newtheorem{definition}[theorem]{\bf Definition}
\newtheorem{conjecture}[theorem]{\bf Conjecture}
\newtheorem{remark}[theorem]{\it Remark}
%%%%%%%%%%%%%%%%%%%%%%%%%%%%%%%%%%%%%%%%%%%%%%%%%%%%%%%%%%%%%%%%%%%%%%%%%%%%%

\newcommand{\FF}{\mathcal{F}}
\newcommand{\Hop}{\mathcal{H}}
\newcommand{\DD}{\mathcal{D}}
\newcommand{\SN}{\mathcal{S}_N}
\newcommand*\dd{\mathrm{d}}
\newcommand*\ii{\mathrm{i}}
%%%%%%%%%%%%%%%%%%%%%%%%%%%%%%%%%%%%%%%%%%%%%%%%%%%%%%%%%%%%%%%%%%%%%%%%%%%%%%%%%%%%%%



% make space for the rule
\fancypagestyle{plain}{
	\fancyhead{} %get rid of the headers on plain pages
	\renewcommand{\headrulewidth}{0pt} % and the line
}

%
% Set header height to 30pt to make room for two rows in the header
% (useful if you have very long chapter names)
%
%\headheight 30pt


%
% Include user-defined macros
%
%====================================================
%-------------Macro definitions go here--------------
%====================================================

%
% Differentials
%
\newcommand{\tdiff}[2]{\ensuremath{\frac{d#2}{d{#1}}}}
\newcommand{\tdifforder}[3]{\ensuremath{\frac{d^{#2}#3}{d{#1}^{#2}}}}
\newcommand{\pdiff}[2]{\ensuremath{\frac{\partial#2 }{\partial#1}}}
\newcommand{\pdifforder}[3]{\ensuremath{\frac{\partial^{#2}#3}{\partial{#1}^{#2}}}}

%
% Linear algebra
%
\renewcommand{\vec}[1]{\ensuremath{\mathbf{#1}}}
\newcommand{\mat}[1]{\ensuremath{\mathbf{#1}}}
\newcommand{\tildemat}[1]{\ensuremath{\widetilde{\mat{#1}}}}


%
% Misc macros
%
\newcommand{\eref}[1]{~(\ref{#1})}
\renewcommand{\equiv}[0]{\ensuremath{:=}}
\newcommand{\etal}{\textit{et al. }}
\newcommand{\paperheader}[2]{\noindent\textbf{Paper #1}: \textit{#2}\\}
\newcommand{\paperitem}[3]{\noindent\textbf{Paper #1}: \textit{#2}\vspace{1em}\\\noindent #3\vspace{2em}}
\newcommand{\tfinal}{\ensuremath{T_{\text{f}}}}
\newcommand{\papernum}[1]{\textbf{#1}}


%
% Theorem environments
%
\newtheorem{definition}{Definition}


%
% Counters
%
\newcounter{ct}

%
% For loop to include papers
%
\newcommand{\includepaperpages}[2]
{
	\forloop{ct}{0}{\value{ct} < #2}
	{
		\begin{figure}[ht]
			\includegraphics[width=.99\textwidth]{#1\the\value{ct}.ps}
		\end{figure}
		\clearpage
	}
}


%
% TeX is very proud of its hyphenation engine, to the point where it 
% hyphenates _everything_ just to show off.  Increasing penalty for 
% hyphenation, and increase tolerance for underfull boxes can make 
% the text easier to read, at the expense of making the text less aligned
% to the right margin
%
\hyphenpenalty=10
\tolerance=1000


% Redefine include command -> input command
%\renewcommand{\include}{\input}

% Include only these chapters
%\includeonly{introduction}


%
% Set up the frontpage, with author, title, etc
%
\title{
{\fontsize{28}{30}\usefont{OT1}{phv}{bc}{n}\selectfont Numerical methods for nonlinear nonlocal water wave models}
	\author{
	\textbf{Daulet Moldabayev}\vspace{3cm}\\
		\includegraphics[width=74mm]{figurer/uglo}\vspace{1em}\\
		Dissertation for the degree of Philosophiae Doctor (PhD)\vspace{3.5em}\\
		Department of Mathematics \\
		University of Bergen \\
		Norway
	}
	\date{April 2017}
}


%====================================================
%------------------ BEGIN DOCUMENT ------------------
%====================================================
\begin{document}
%Cause all references in bibtex file to appear in the 'References' section, even
%if they are not explicitly cite'ed in the document
%\nocite{*}

%Set the fontsize and baselineskip, if something other than 10 or 12 is required
\ifDownscaledFinalDoc
	\fontsize{\TextSize}{\BaseLineSkip}
	\selectfont
\fi

%Double line spacing for draft
\ifDraft
	\doublespacing
\fi


%--------------------------------------------------------------------------------------------------
% FRONT MATTER
%--------------------------------------------------------------------------------------------------
\maketitle
\frontmatter
%\chapter{Scientific environment}

This is where you write the name of the subject/professional milieu involved in the study, faculty(faculties)/institute(s)/centre(s)/research groups/school of research.

Contractual joint venture partners should also be mentioned here. If these are represented by graphic symbols (logo etc.) the it should be on this page.

NB! Most organisations have guidelines concerning the use of their logo. These guidelines regulate who can use the logo, when the logo can be used and how. This is quite relevant when a logo is used together with information connected to other institutions/organisations. (A typical example would be rules concerning space between other written or graphic content.)

Remember to find out which guidelines apply and ensure that they are followed!
Most organisations will be able to provide a digital version of their logo in a suitable format.

%\chapter{Preface}

This dissertation is submitted as a partial fulfillment of the requirements for the degree Doctor of Philosophy (PhD) at the Department of Mathematics, University of Bergen. The research was
supported by the Research Council of Norway on grant no. 213474/F20. 




The advisory committee has consisted of Henrik Kalisch (University of Bergen, Bergen, Norway),  Dmitry Pelinovsky (McMaster University, Hamilton, Ontario, Canada) and Angel Dur\'an (Universidad de Valladolid, Valladolid, Spain).


%\chapter{Acknowledgements}

First of all I would like to thank Henrik Kalisch and Magnus Sv\"ard who carried out the selection process of candidates for the PhD position. I am sincerely thankful to Professor Henrik Kalisch, my main supervisor, for sharing knowledge and guiding me, for his enthusiasm, support and encouragement. The thesis would not have been possible without you.


In the course of the doctorate studies, I have had fruitful collaboration with Olivier Verdier, Evgueni Dinvay, Denys Dutykh, Ben Segal and others. I would like to thank them for productive discussions, advices on research questions and some inspiring ideas we came up with.


As a PhD Candidate I have been lucky to meet my colleagues and fellow students at the Department of Mathematics. They have always provided a friendly and supportive environment. I appreciate everything we have experienced during the time of our studies.


I would like to express my gratitude to my family for their continuous support and encouragement over the years of my studies. I am also thankful for all of my teachers and professors from schools and universities I studied at. 


I am heartily grateful for my wife Assem. You have always been there for me. Without your love, support and understanding I would not have been able to go that far. 
%\chapter{Abstract}

Nonlinear water wave models 
%\chapter{Outline}



This thesis is organised in the following way. Part I contains general theoretical background on nonlinear wave models as well as description of methods used to solve the equations involved. Some properties of solutions to the equations and a summary of results are also given in the first part. Part 2 consists of the research papers that present scientific results in detail. 



\section*{List of reseach papers included in Part II}


\textbf{Paper A:}

\noindent Moldabayev, D., Kalisch, H., Dutykh, D.: \textit{The Whitham equation as a model for surface water waves}, Phys. D~ \textbf{309}, 99–107 (2015),  

\noindent http://dx.doi.org/10.1016/j.physd.2015.07.010.

\vspace{1em}

\noindent \textbf{Paper B:}

\noindent Dinvay, E., Moldabayev, D., Dutykh, D., Kalisch, H.:  \textit{The Whitham equation with surface tension}, Nonlinear Dynamics, 1--14 (2017),  

\noindent http://dx.doi.org/10.1007/s11071-016-3299-7.

\vspace{1em}

\noindent \textbf{Paper C:}

\noindent Kalisch, H., Moldabayev, D., Verdier, O.: \textit{A numerical study of nonlinear dispersive wave models with SpecTraVVave}, Electronic Journal of Differential Equations~ \textbf{2017}, No. 62, pp. 1--23 (2017), http://ejde.math.txstate.edu or http://ejde.math.unt.edu.

\vspace{1em}

\noindent \textbf{Paper D:}

\noindent Benjamin Segal, Daulet Moldabayev, Henrik Kalisch, Bernard Deconinck: \textit{Explicit solutions for a long-wave model with constant vorticity}, submitted to European Journal of Mechanics - B/Fluids.

%\chapter{List of papers}

\begin{enumerate}

\item Daulet Moldabayev, Henrik Kalisch, Denys Dutyh, \textit{The Whitham equation as a model for surface water waves}, Physica D \textbf{309}, 2015.

\end{enumerate}

\tableofcontents
%\listoffigures
%\lstlistoflistings

%--------------------------------------------------------------------------------------------------
% MAIN MATTER
%--------------------------------------------------------------------------------------------------
\mainmatter

%
% Main chapters
%
\part{Background}
%\part{Background}
%%%%%%%%%%%%%%%%%%%%%%%%%%%%%%%%%%%%%%%%%%%%%%%%%%%%%%%%%%%%%%%%%%%%%%%%%%%%%%%%%%%%%%%%%%%
\chapter{Introduction}
%%%%%%%%%%%%%%%%%%%%%%%%%%%%%%%%%%%%%%%%%%%%%%%%%%%%%%%%%%%%%%%%%%%%%%%%%%%%%%%%%%%%%%%%%%%
This is the introduction~\cite{Dinvay2017}\ldots


%%%%%%%%%%%%%%%%%%%%%%%%%%%%%%%%%%%%%%%%%%%%%%%%%%%%%%%%%%%%%%%%%%%%%%%%%%%%%%%%%%%%%%%%%%%%
\section{Motivation}
%%%%%%%%%%%%%%%%%%%%%%%%%%%%%%%%%%%%%%%%%%%%%%%%%%%%%%%%%%%%%%%%%%%%%%%%%%%%%%%%%%%%%%%%%%%

%%%%%%%%%%%%%%%%%%%%%%%%%%%%%%%%%%%%%%%%%%%%%%%%%%%%%%%%%%%%%%%%%%%%%%%%%%%%%%%%%%%%%%%%%%%
%\section{Motivation}
%%%%%%%%%%%%%%%%%%%%%%%%%%%%%%%%%%%%%%%%%%%%%%%%%%%%%%%%%%%%%%%%%%%%%%%%%%%%%%%%%%%%%%%%%%%


%%%%%%%%%%%%%%%%%%%%%%%%%%%%%%%%%%%%%%%%%%%%%%%%%%%%%%%%%%%%%%%%%%%%%%%%%%%%%%%%%%%%%%%%%%%%
\chapter{Theory on water-wave models}
%%%%%%%%%%%%%%%%%%%%%%%%%%%%%%%%%%%%%%%%%%%%%%%%%%%%%%%%%%%%%%%%%%%%%%%%%%%%%%%%%%%%%%%%%%%



%%%%%%%%%%%%%%%%%%%%%%%%%%%%%%%%%%%%%%%%%%%%%%%%%%%%%%%%%%%%%%%%%%%%%%%%%%%%%%%%%%%%%%%%%%%
\section{Derivation of models and equations}
%%%%%%%%%%%%%%%%%%%%%%%%%%%%%%%%%%%%%%%%%%%%%%%%%%%%%%%%%%%%%%%%%%%%%%%%%%%%%%%%%%%%%%%%%%%






%%%%%%%%%%%%%%%%%%%%%%%%%%%%%%%%%%%%%%%%%%%%%%%%%%%%%%%%%%%%%%%%%%%%%%%%%%%%%%%%%%%%%%%%%%%
\section{Hamiltonian theory}
%%%%%%%%%%%%%%%%%%%%%%%%%%%%%%%%%%%%%%%%%%%%%%%%%%%%%%%%%%%%%%%%%%%%%%%%%%%%%%%%%%%%%%%%%%%





%%%%%%%%%%%%%%%%%%%%%%%%%%%%%%%%%%%%%%%%%%%%%%%%%%%%%%%%%%%%%%%%%%%%%%%%%%%%%%%%%%%%%%%%%%%
\section{Solutions and their stability}
%%%%%%%%%%%%%%%%%%%%%%%%%%%%%%%%%%%%%%%%%%%%%%%%%%%%%%%%%%%%%%%%%%%%%%%%%%%%%%%%%%%%%%%%%%%




%%%%%%%%%%%%%%%%%%%%%%%%%%%%%%%%%%%%%%%%%%%%%%%%%%%%%%%%%%%%%%%%%%%%%%%%%%%%%%%%%%%%%%%%%%%
\section{Bifurcation curves and their connection to stability}
%%%%%%%%%%%%%%%%%%%%%%%%%%%%%%%%%%%%%%%%%%%%%%%%%%%%%%%%%%%%%%%%%%%%%%%%%%%%%%%%%%%%%%%%%%%

%%%%%%%%%%%%%%%%%%%%%%%%%%%%%%%%%%%%%%%%%%%%%%%%%%%%%%%%%%%%%%%%%%%%%%%%%%%%%%%%%%%%%%%%%%%
%\chapter{Theory on water-wave models}
%%%%%%%%%%%%%%%%%%%%%%%%%%%%%%%%%%%%%%%%%%%%%%%%%%%%%%%%%%%%%%%%%%%%%%%%%%%%%%%%%%%%%%%%%%%



%%%%%%%%%%%%%%%%%%%%%%%%%%%%%%%%%%%%%%%%%%%%%%%%%%%%%%%%%%%%%%%%%%%%%%%%%%%%%%%%%%%%%%%%%%%
%\section{Derivation of models and equations}
%%%%%%%%%%%%%%%%%%%%%%%%%%%%%%%%%%%%%%%%%%%%%%%%%%%%%%%%%%%%%%%%%%%%%%%%%%%%%%%%%%%%%%%%%%%






%%%%%%%%%%%%%%%%%%%%%%%%%%%%%%%%%%%%%%%%%%%%%%%%%%%%%%%%%%%%%%%%%%%%%%%%%%%%%%%%%%%%%%%%%%%
%\section{Hamiltonian theory}
%%%%%%%%%%%%%%%%%%%%%%%%%%%%%%%%%%%%%%%%%%%%%%%%%%%%%%%%%%%%%%%%%%%%%%%%%%%%%%%%%%%%%%%%%%%





%%%%%%%%%%%%%%%%%%%%%%%%%%%%%%%%%%%%%%%%%%%%%%%%%%%%%%%%%%%%%%%%%%%%%%%%%%%%%%%%%%%%%%%%%%%
%\section{Solutions and their stability}
%%%%%%%%%%%%%%%%%%%%%%%%%%%%%%%%%%%%%%%%%%%%%%%%%%%%%%%%%%%%%%%%%%%%%%%%%%%%%%%%%%%%%%%%%%%




%%%%%%%%%%%%%%%%%%%%%%%%%%%%%%%%%%%%%%%%%%%%%%%%%%%%%%%%%%%%%%%%%%%%%%%%%%%%%%%%%%%%%%%%%%%
%\section{Bifurcation curves and their connection to stability}
%%%%%%%%%%%%%%%%%%%%%%%%%%%%%%%%%%%%%%%%%%%%%%%%%%%%%%%%%%%%%%%%%%%%%%%%%%%%%%%%%%%%%%%%%%%


%%%%%%%%%%%%%%%%%%%%%%%%%%%%%%%%%%%%%%%%%%%%%%%%%%%%%%%%%%%%%%%%%%%%%%%%%%%%%%%%%%%%%%%%%%%%
\chapter{Numerics for water-wave models}
%%%%%%%%%%%%%%%%%%%%%%%%%%%%%%%%%%%%%%%%%%%%%%%%%%%%%%%%%%%%%%%%%%%%%%%%%%%%%%%%%%%%%%%%%%%



%%%%%%%%%%%%%%%%%%%%%%%%%%%%%%%%%%%%%%%%%%%%%%%%%%%%%%%%%%%%%%%%%%%%%%%%%%%%%%%%%%%%%%%%%%%
\section{Fixed point methods}
%%%%%%%%%%%%%%%%%%%%%%%%%%%%%%%%%%%%%%%%%%%%%%%%%%%%%%%%%%%%%%%%%%%%%%%%%%%%%%%%%%%%%%%%%%%

%%%%%%%%%%%%%%%%%%%%%%%%%%%%%%%%%%%%%%%%%%%%%%%%%%%%%%%%%%%%%%%%%%%%%%%%%%%%%%%%%%%%%%%%%%%
%\chapter{Numerics for water-wave models}
%%%%%%%%%%%%%%%%%%%%%%%%%%%%%%%%%%%%%%%%%%%%%%%%%%%%%%%%%%%%%%%%%%%%%%%%%%%%%%%%%%%%%%%%%%%



%%%%%%%%%%%%%%%%%%%%%%%%%%%%%%%%%%%%%%%%%%%%%%%%%%%%%%%%%%%%%%%%%%%%%%%%%%%%%%%%%%%%%%%%%%%
%\section{Fixed point methods}
%%%%%%%%%%%%%%%%%%%%%%%%%%%%%%%%%%%%%%%%%%%%%%%%%%%%%%%%%%%%%%%%%%%%%%%%%%%%%%%%%%%%%%%%%%%



%%%%%%%%%%%%%%%%%%%%%%%%%%%%%%%%%%%%%%%%%%%%%%%%%%%%%%%%%%%%%%%%%%%%%%%%%%%%%%%%%%%%%%%%%%%%
\chapter{Pressure reconstruction theory}
%%%%%%%%%%%%%%%%%%%%%%%%%%%%%%%%%%%%%%%%%%%%%%%%%%%%%%%%%%%%%%%%%%%%%%%%%%%%%%%%%%%%%%%%%%%

%%%%%%%%%%%%%%%%%%%%%%%%%%%%%%%%%%%%%%%%%%%%%%%%%%%%%%%%%%%%%%%%%%%%%%%%%%%%%%%%%%%%%%%%%%%
%\chapter{Pressure reconstruction theory}
%%%%%%%%%%%%%%%%%%%%%%%%%%%%%%%%%%%%%%%%%%%%%%%%%%%%%%%%%%%%%%%%%%%%%%%%%%%%%%%%%%%%%%%%%%%





%%%%%%%%%%%%%%%%%%%%%%%%%%%%%%%%%%%%%%%%%%%%%%%%%%%%%%%%%%%%%%%%%%%%%%%%%%%%%%%%%%%%%%%%%%%
\chapter{Summary of results}
%%%%%%%%%%%%%%%%%%%%%%%%%%%%%%%%%%%%%%%%%%%%%%%%%%%%%%%%%%%%%%%%%%%%%%%%%%%%%%%%%%%%%%%%%%%

This chapter provides an overview of the results achieved in the course of research work. 
Detailed 

%%%%%%%%%%%%%%%%%%%%%%%%%%%%%%%%%%%%%%%%%%%%%%%%%%%%%%%%%%%%%%%%%%%%%%%%%%%%%%%%%%%%%%%%%%%
\section{The Whitham equation as a model for surface water waves}
%%%%%%%%%%%%%%%%%%%%%%%%%%%%%%%%%%%%%%%%%%%%%%%%%%%%%%%%%%%%%%%%%%%%%%%%%%%%%%%%%%%%%%%%%%%

The Whitham equation was proposed as an alternate model equation for the simplified description of unidirectional 
wave motion at the surface of an inviscid fluid. As the Whitham equation incorporates the full
linear dispersion relation of the water wave problem, it is thought to provide a more faithful description
of shorter waves of small amplitude than traditional long wave models such as the KdV equation.
In this work, we identify a scaling regime in which the Whitham equation can be derived from the
Hamiltonian theory of surface water waves. A Hamiltonian system of Whitham type allowing for two-
way wave propagation is also derived. The Whitham equation is integrated numerically, and it is shown
that the equation gives a close approximation of inviscid free surface dynamics as described by the
Euler equations. The performance of the Whitham equation as a model for free surface dynamics is also
compared to different free surface models: the KdV equation, the BBM equation, and the Padé (2,2) model.
It is found that in a wide parameter range of amplitudes and wavelengths, the Whitham equation performs
on par with or better than the three considered models.


In its simplest form, the water-wave problem concerns the
flow of an incompressible inviscid fluid with a free surface over
a horizontal impenetrable bed. In this situation, the fluid flow
is described by the Euler equations with appropriate boundary
conditions, and the dynamics of the free surface are of particular
interest in the solution of this problem.
There are a number of model equations which allow the
approximate description of the evolution of the free surface
without having to provide a complete solution of the fluid flow
below the surface. In the present contribution, interest is focused

on the derivation and evaluation of a nonlocal water-wave model
known as the Whitham equation. The equation is written as

\begin{equation}
\eta_t + \frac{3}{2} \frac{c_0}{h_0} \eta \eta_x + K_{h_0} \ast \eta_x = 0, 
\label{whitham-equation}
\end{equation}
where the convolution kernel $K_{h_0}$ is given in terms of the Fourier
transform by

\begin{equation}
\mathcal{F}K_{h_0} (\xi) =  \sqrt{\frac{g \tanh(h_0\xi)}{\xi}}.
\label{whitham-kernel}
\end{equation}
$g$ is the gravitational acceleration, $h_0$ is the undisturbed depth of
the fluid, and $c_0 = \sqrt{gh_0}$ is the corresponding long-wave speed.
The convolution can be thought of as a Fourier multiplier operator,
and \eqref{whitham-kernel} represents the Fourier symbol of the operator.
The Whitham equation was proposed by Whitham \cite{Whitham1967} as an
alternative to the well known Korteweg–de Vries (KdV) equation

\begin{equation}
\eta_t + c_0 \eta_x + \frac{3}{2} \frac{c_0}{h_0} \eta \eta_x + \frac{1}{6} c_0 h_0^2 \eta_{xxx} = 0.
\label{kdv-equation}
\end{equation}


The validity of the KdV equation as a model for surface water waves
can be described as follows. Suppose a wave field with a prominent
amplitude $a$ and characteristic wavelength $l$ is to be studied. The
KdV equation is known to produce a good approximation of the
evolution of the waves if the amplitude of the waves is small and
the wavelength is large when compared to the undisturbed depth,
and if in addition, the two non-dimensional quantities $a / h_0$ and
$h_0^2 / l^2$ are of similar size. The latter requirement can be written in
terms of the Stokes number as

\begin{equation}
\mathcal{S} = \frac{a l^2}{h_0^3} \sim 1. 
\label{stockes-number}
\end{equation}
While the KdV equation is a good model for surface waves if $\mathcal{S} \sim 1$, 
one notorious problem with the KdV equation is that it does
not model accurately the dynamics of shorter waves. Recognizing
this shortcoming of the KdV equation, Whitham proposed to use
the same nonlinearity as the KdV equation, but couple it with a
linear term which mimics the linear dispersion relation of the full
water-wave problem. Thus, at least in theory, the Whitham
equation can be expected to yield a description of the dynamics
of shorter waves which is closer to the solutions of the more
fundamental Euler equations which govern the flow.

The Whitham equation has been studied from a number of
vantage points during recent years. In particular, the existence of
traveling and solitary waves has been established in \citep{Ehrnstrom2012, Ehrnstrom2009}.
Well
posedness of a similar equation was investigated in \citetext{\citealp[]{Lannes2013}-\citealp[]{Klein2015}}, and a
model with variable depth has been studied numerically in \cite{Minzoni2013}.
Moreover, it has been shown in \citep{Hur2015, Sanford2014} that periodic solutions of
the Whitham equation feature modulational instability for short
enough waves in a similar way as small-amplitude periodic wave
solutions of the water-wave problem. However, even though the
equation is routinely mentioned in texts on nonlinear waves \citep{Drazin89, Whitham74}, it appears that the performance of the Whitham equation in
the description of surface water waves has not been investigated
so far.
The purpose of the present article is to give an asymptotic
derivation of the Whitham equation as a model for surface water
waves, and to confirm Whitham’s expectation that the equation is
a fair model for the description of time-dependent surface water
waves. For the purpose of the derivation, we introduce an expo-
nential scaling regime in which the Whitham equation can be de-
rived asymptotically from an approximate Hamiltonian principle
for surface water waves. In order to motivate the use of this scaling,
note that the KdV equation has the property that wide classes of
initial data decompose into a number of solitary waves and small-
amplitude dispersive residue \cite{Ablowitz81}. For the KdV equations, solitary-
wave solutions are known in closed form, and are given by
\begin{equation}
\eta = \frac{a}{h_0}\mbox{sech}^2\left(\sqrt{\frac{3a}{4h_0^3}(x - ct)}\right)
\label{kdv-soliton}
\end{equation}
for a certain wave celerity $c$. These waves clearly comply with
the amplitude–wavelength relation $a / h_0 \sim h_0^2 / l^2$ which was
mentioned above. It appears that the Whitham equation – as indeed do many other nonlinear dispersive equations – also has
the property that broad classes of initial data rapidly decompose
into ordered trains of solitary waves (see Fig. 1). Quantifying the
amplitude–wavelength relation for these solitary waves yields an
asymptotic regime which is expected to be relevant to the validity
of the Whitham equation as a water wave model.

%%%%%%%%%%%%%%%%%%%%%%%%%%%%%%%%%%%%%%%%%%%%%%%%%%%%%%%%%%%%%%%%%%%%%%%%%%%%%%%%%%%%%%%%%%%
\section{The Whitham equation with surface tension}
%%%%%%%%%%%%%%%%%%%%%%%%%%%%%%%%%%%%%%%%%%%%%%%%%%%%%%%%%%%%%%%%%%%%%%%%%%%%%%%%%%%%%%%%%%%






%%%%%%%%%%%%%%%%%%%%%%%%%%%%%%%%%%%%%%%%%%%%%%%%%%%%%%%%%%%%%%%%%%%%%%%%%%%%%%%%%%%%%%%%%%%
\section{A numerical study of nonlinear dispersive wave models with SpecTraVVave}
%%%%%%%%%%%%%%%%%%%%%%%%%%%%%%%%%%%%%%%%%%%%%%%%%%%%%%%%%%%%%%%%%%%%%%%%%%%%%%%%%%%%%%%%%%%





%%%%%%%%%%%%%%%%%%%%%%%%%%%%%%%%%%%%%%%%%%%%%%%%%%%%%%%%%%%%%%%%%%%%%%%%%%%%%%%%%%%%%%%%%%%
\section{Explicit solutions for a long-wave model with constant vorticity}
%%%%%%%%%%%%%%%%%%%%%%%%%%%%%%%%%%%%%%%%%%%%%%%%%%%%%%%%%%%%%%%%%%%%%%%%%%%%%%%%%%%%%%%%%%%

%%%%%%%%%%%%%%%%%%%%%%%%%%%%%%%%%%%%%%%%%%%%%%%%%%%%%%%%%%%%%%%%%%%%%%%%%%%%%%%%%%%%%%%%%%%
%S\chapter{Summary of results}
%%%%%%%%%%%%%%%%%%%%%%%%%%%%%%%%%%%%%%%%%%%%%%%%%%%%%%%%%%%%%%%%%%%%%%%%%%%%%%%%%%%%%%%%%%%





%%%%%%%%%%%%%%%%%%%%%%%%%%%%%%%%%%%%%%%%%%%%%%%%%%%%%%%%%%%%%%%%%%%%%%%%%%%%%%%%%%%%%%%%%%%
%\section{The Whitham equation as a model for surface water waves}
%%%%%%%%%%%%%%%%%%%%%%%%%%%%%%%%%%%%%%%%%%%%%%%%%%%%%%%%%%%%%%%%%%%%%%%%%%%%%%%%%%%%%%%%%%%







%%%%%%%%%%%%%%%%%%%%%%%%%%%%%%%%%%%%%%%%%%%%%%%%%%%%%%%%%%%%%%%%%%%%%%%%%%%%%%%%%%%%%%%%%%%
%\section{The Whitham equation with surface tension}
%%%%%%%%%%%%%%%%%%%%%%%%%%%%%%%%%%%%%%%%%%%%%%%%%%%%%%%%%%%%%%%%%%%%%%%%%%%%%%%%%%%%%%%%%%%






%%%%%%%%%%%%%%%%%%%%%%%%%%%%%%%%%%%%%%%%%%%%%%%%%%%%%%%%%%%%%%%%%%%%%%%%%%%%%%%%%%%%%%%%%%%
%\section{A numerical study of nonlinear dispersive wave models with SpecTraVVave}
%%%%%%%%%%%%%%%%%%%%%%%%%%%%%%%%%%%%%%%%%%%%%%%%%%%%%%%%%%%%%%%%%%%%%%%%%%%%%%%%%%%%%%%%%%%





%%%%%%%%%%%%%%%%%%%%%%%%%%%%%%%%%%%%%%%%%%%%%%%%%%%%%%%%%%%%%%%%%%%%%%%%%%%%%%%%%%%%%%%%%%%
%\section{Explicit solutions for a long-wave model with constant vorticity}
%%%%%%%%%%%%%%%%%%%%%%%%%%%%%%%%%%%%%%%%%%%%%%%%%%%%%%%%%%%%%%%%%%%%%%%%%%%%%%%%%%%%%%%%%%%
         % Introduction
\chapter{Introduction to the papers}

%
% Paper I
%
\paperitem{A}
{The Whitham equation as a model for surface water waves}
{In this paper\ldots}

  % Introduction to the papers

%
% Appendices
%
%\appendix
%\chapter{Some appendix}

             % Appendix A: Something


%--------------------------------------------------------------------------------------------------
% BACK MATTER
%--------------------------------------------------------------------------------------------------
\backmatter

%
% The bibliography
%

%
% Physics-style bibliography
%
%\bibliographystyle{iopart-num}

%
% Modified mathematics bibliography. Uses acm.bst
%
\addcontentsline{toc}{chapter}{Bibliography}
\bibliographystyle{acm}

\bibliography{thesis}

\part{Scientific results}
%\pagenumbering{gobble}
%%\part{Scientific results}


%--------------------------------------------------------------------------------------------------
% Paper A
%--------------------------------------------------------------------------------------------------
\chapter*{Paper A}

\addcontentsline{toc}{chapter}{Paper A The Whitham equation as a model for surface water waves}
\section*{The Whitham equation as a model for surface water waves}

\noindent Daulet Moldabayev, Henrik Kalisch, Denys Dutyh \\

\noindent \textit{Physica} D \textbf{309}, (2015)

\noindent http://dx.doi.org/10.1016/j.physd.2015.07.010
\cleardoublepage

% This will include the bursted pdf. Remember to update the number of pages
% (last parameter) to the actual number of pages
\includepaperpages{../papers/paper-1/page_}{9}

%--------------------------------------------------------------------------------------------------
% Paper B
%--------------------------------------------------------------------------------------------------

\chapter*{Paper B}

\addcontentsline{toc}{chapter}{Paper B The Whitham equation with surface tension}
\section*{The Whitham equation with surface tension}

\noindent Evgueni Dinvay, Daulet Moldabayev, Denys Dutykh, Henrik Kalisch\\

\noindent \textit{Nonlinear Dynamics}, 1--14 (2017)

\noindent http://dx.doi.org/10.1007/s11071-016-3299-7

\cleardoublepage

% This will include the bursted pdf. Remember to update the number of pages
% (last parameter) to the actual number of pages
\includepaperpages{../papers/paper-2/page_}{14}

%--------------------------------------------------------------------------------------------------
% Paper C
%--------------------------------------------------------------------------------------------------

\chapter*{Paper C}

\addcontentsline{toc}{chapter}{Paper C A numerical study of nonlinear dispersive wave models with SpecTraVVave}
\section*{A numerical study of nonlinear dispersive wave models with SpecTraVVave}

\noindent Henrik Kalisch, Daulet Moldabayev, Olivier Verdier\\

\noindent \textit{Electronic Journal of Differential Equations} \textbf{2017},  No. 62, pp. 1--23, (2017)

\noindent http://ejde.math.txstate.edu or http://ejde.math.unt.edu

%\noindent \textit{Physical Review} B, \textbf{76}, 035303 (2007)
\cleardoublepage

% This will include the bursted pdf. Remember to update the number of pages
% (last parameter) to the actual number of pages
\includepaperpages{../papers/paper-3/page_}{23}

%--------------------------------------------------------------------------------------------------
% Paper D
%--------------------------------------------------------------------------------------------------

\chapter*{Paper D}

\addcontentsline{toc}{chapter}{Paper D Explicit solutions for a long-wave model with constant vorticity}
\section*{Explicit solutions for a long-wave model with constant vorticity}

\noindent Benjamin Segal, Daulet Moldabayev, Henrik Kalisch, Bernard Deconinck\\

%\noindent \textit{Physical Review} B, \textbf{76}, 035303 (2007)
\cleardoublepage

% This will include the bursted pdf. Remember to update the number of pages
% (last parameter) to the actual number of pages
\includepaperpages{../papers/paper-4/page_}{19}
               % The papers, with cover pages


\end{document}

